% =====================================================================
%  WEB-BASED SUPPORTING MATERIALS for the Series C paper.
%  Series C explicitly invites proofs into the supplement; they are here,
%  not in the main text.
%
%  BUILD: pdflatex supplement ; pdflatex supplement
%
%  NUMBERS: every value is registered in ../paper1_speckit/08_RESULTS_MANIFEST.md.
%  The manifest section is named in a comment at each use.
% =====================================================================
\documentclass[11pt,a4paper]{article}
\usepackage[T1]{fontenc}
\usepackage{mathptmx}
\usepackage[margin=2.6cm]{geometry}
\usepackage{amsmath,amssymb,amsthm}
\usepackage{booktabs}
\usepackage{graphicx}
\usepackage[authoryear,round]{natbib}
\usepackage{hyperref}
\hypersetup{colorlinks=true,linkcolor=black,citecolor=black,urlcolor=blue}

\DeclareMathOperator*{\argmax}{arg\,max}
\DeclareMathOperator{\var}{var}
\DeclareMathOperator{\sd}{sd}
\DeclareMathOperator{\tr}{tr}
\newcommand{\bb}{\boldsymbol{\beta}}
\newcommand{\bbh}{\hat{\boldsymbol{\beta}}}
\newcommand{\bbt}{\tilde{\boldsymbol{\beta}}}
\newcommand{\by}{\boldsymbol{y}}
\newcommand{\bpi}{\boldsymbol{\pi}}
\newcommand{\br}{\boldsymbol{r}}
\newcommand{\bmu}{\boldsymbol{\mu}}

\theoremstyle{plain}
\newtheorem{proposition}{Proposition}
\newtheorem{lemma}{Lemma}
\theoremstyle{remark}
\newtheorem{remark}{Remark}

\renewcommand{\thesection}{S\arabic{section}}
\renewcommand{\thetable}{S\arabic{table}}
\renewcommand{\thefigure}{S\arabic{figure}}
\renewcommand{\theequation}{S\arabic{equation}}

\title{\bfseries Supplementary material for\\
``Shrinkage invalidates the Hosmer--Lemeshow test: goodness of fit for penalized
logistic regression, with an application to glaucoma diagnosis''}
\author{Ebrahim Khaled Ebrahim}
\date{}

\begin{document}
\maketitle

\noindent
Numbering of propositions matches the main text. Section~\ref{sec:proofs} gives the
regularity conditions and the proofs. Section~\ref{sec:failures} catalogues the
neighbouring constructions that fail, with the size of each failure, because the
question ``why not the obvious thing?'' has four different answers.
Section~\ref{sec:designs} specifies the simulation designs.
Section~\ref{sec:grids} gives the full grids and the sensitivity analyses.
Section~\ref{sec:repro} lists the scripts and seeds.

\tableofcontents


% =====================================================================
\section{Regularity conditions and proofs}\label{sec:proofs}
% =====================================================================

\subsection{Conditions}\label{sec:conditions}

Throughout, $y_i \in \{0,1\}$ are independent given $\boldsymbol{x}_i$, and
$\pi_i(\bb) = \{1+\exp(-\boldsymbol{x}_i^{\top}\bb)\}^{-1}$. Write
$W = \mathrm{diag}\{\hat\pi_i(1-\hat\pi_i)\}$, $F = X^{\top}WX$, $K = \lambda D$ with
$D = \mathrm{diag}(0,1,\dots,1)$, and $M = F+K$. We assume:

(C1) the true parameter $\bb_0$ lies in the interior of a compact parameter set, and
$\pi_i(\bb_0) \in [\epsilon, 1-\epsilon]$ for some $\epsilon > 0$ uniformly in $i$;

(C2) $n^{-1}F \to \Sigma$, positive definite, and $\lambda$ is non-random with
$\lambda = o(n)$, so that $n^{-1}M \to \Sigma$ and the shrinkage displacement
$M^{-1}K\bb = o(1)$;

(C3) the grouping matrix $C$ has $G$ fixed, each group receiving a fixed positive
fraction of the sample, and $V = CWC^{\top}$ is positive definite;

(C3$'$) the group boundaries converge in probability to the population quantiles of the
linear predictor at rate $O_p(n^{-1/2})$, with limiting cell probabilities bounded away
from zero. This is an assumption, not something we discharge: in practice the groups are
equal-frequency cells of the \emph{fitted} index, so $C$ is random. \citet{moore1975}
establish the corresponding statement for the unpenalized case; it is untested at
$p/n = 0.25$;

(C4) $p$ is fixed as $n \to \infty$.

Condition (C4) is what confines the results of this section to the fixed-dimension
regime. It is not satisfied by design~B, where $p/n = 0.25$, and this is exactly why the
validity we claim there rests on prepivoting and on the simulation evidence of
Section~\ref{sec:grids}, not on Propositions 1 and 2. We are explicit about this in the
discussion of the main text.

\begin{remark}
Conditions (C1)--(C3) are those of \citet{moore1975} with one deletion and one addition.
Moore and Spruill require the estimator to be $\sqrt{n}$-consistent for $\bb_0$ and to
satisfy an orthogonality condition on the estimating function. A penalized estimator
satisfies neither when $\lambda$ grows with $n$: it is $\sqrt{n}$-consistent for a
shrunken pseudo-parameter, not for $\bb_0$. Proposition~\ref{prop:law} is what replaces
their conclusion once those two conditions are dropped.
\end{remark}

\subsection{Proof of Proposition 1}\label{sec:proof1}

\begin{proposition}\label{prop:law}
Under (C1)--(C4),
$\bbh - \bb = M^{-1}\{X^{\top}(\by-\bpi) - K\bb\} + o_p(1)$ and
$\br \mathrel{\dot\sim} \mathcal{N}(\bmu_K, \Omega_K)$ with
\begin{align*}
  \Omega_K &= I_G - UM^{-1}U^{\top} - UM^{-1}KM^{-1}U^{\top}
            = I_G - UM^{-1}(F+2K)M^{-1}U^{\top}, \\
  \bmu_K   &= UM^{-1}K\bb + \text{(second-order remainder)}.
\end{align*}
\end{proposition}

\begin{proof}
The penalized score is $s(\bb) = X^{\top}\{\by - \bpi(\bb)\} - K\bb$, and $\bbh$ solves
$s(\bbh) = \boldsymbol{0}$. Expanding about the true $\bb$ and using
$-\partial s/\partial \bb = F + K = M$,
\[
  \boldsymbol{0} = s(\bbh) = X^{\top}(\by-\bpi) - K\bb - M(\bbh - \bb) + o_p(n^{1/2}),
\]
which on rearranging gives the stated expansion, the remainder being $o_p(1)$ after
scaling by $n^{-1/2}$ under (C2).

For the residuals, expand $\hat\bpi = \bpi(\bbh)$ about $\bpi = \bpi(\bb)$:
$\hat\bpi = \bpi + WX(\bbh-\bb) + o_p(n^{-1/2})$, so that
\[
  \br = V^{-1/2}C(\by - \hat\bpi)
      = V^{-1/2}C(\by-\bpi) - U(\bbh - \bb) + o_p(1),
\]
using $U = V^{-1/2}CWX$. Substituting the expansion for $\bbh - \bb$,
\begin{equation}\label{eq:rlinear}
  \br = \underbrace{\bigl(V^{-1/2}C - UM^{-1}X^{\top}\bigr)}_{=:A}(\by-\bpi)
        \; + \; UM^{-1}K\bb + o_p(1).
\end{equation}
Under $H_0^{\mathrm{str}}$ we have $\mathbb{E}(\by-\bpi) = \boldsymbol{0}$, so
$\mathbb{E}(\br) = UM^{-1}K\bb$, which is $\bmu_K$. We write the remainder in words rather
than as an order symbol because $\|K\bb\| = O(\lambda)$, so $O(\|K\bb\|^{2})$ is not a useful
bound against a leading term of order $n^{-1/2}$. What we can say is measured rather than
proved: at design~A with $\lambda = 100$ the largest discrepancy between the empirical mean
of $\br$ and $\bmu_K$ is $0.1137$ against $\|\bmu_K\| = 3.871$, about $3\%$; and at
$\lambda = 416$ in design~B the second-order term overshoots the empirical mean by roughly a
factor of two, which is why Section~\ref{sec:failures} row~(b) fails.

For the variance, $\var(\by-\bpi) = W$, so $\var(\br) = AWA^{\top}$. Expanding, and using
the three identities $V^{-1/2}CWC^{\top}V^{-1/2} = I_G$, $V^{-1/2}CWX = U$ and
$X^{\top}WX = F$,
\[
  AWA^{\top}
  = I_G - UM^{-1}U^{\top} - UM^{-1}U^{\top} + UM^{-1}FM^{-1}U^{\top}
  = I_G - 2UM^{-1}U^{\top} + UM^{-1}FM^{-1}U^{\top}.
\]
Since $F = M - K$,
$UM^{-1}FM^{-1}U^{\top} = UM^{-1}U^{\top} - UM^{-1}KM^{-1}U^{\top}$, and substituting
gives $\Omega_K = I_G - UM^{-1}U^{\top} - UM^{-1}KM^{-1}U^{\top}$. For the second form,
$UM^{-1}(F+2K)M^{-1}U^{\top} = UM^{-1}(M+K)M^{-1}U^{\top}
= UM^{-1}U^{\top} + UM^{-1}KM^{-1}U^{\top}$. Asymptotic normality follows from the
Lindeberg--Feller theorem applied to the linear form in \eqref{eq:rlinear} under
(C1)--(C3).
\end{proof}

\begin{remark}\label{rem:sign}
Setting $K = 0$ recovers $I_G - UF^{-1}U^{\top}$, the classical
\citet{chernoff1954,moore1975} form. The ordering against that form is the opposite of
what a first glance at Proposition~\ref{prop:law} suggests, and it is worth deriving because the
wrong sign here would misdirect the whole diagnosis. Since $M = F + K$ gives
$F^{-1} - M^{-1} = F^{-1}KM^{-1} = M^{-1}KF^{-1}$, and
$M^{-1}(F+2K)M^{-1} = M^{-1}(M+K)M^{-1} = M^{-1} + M^{-1}KM^{-1}$,
\[
  F^{-1} - M^{-1}(F+2K)M^{-1}
  = \bigl(F^{-1} - M^{-1}\bigr) - M^{-1}KM^{-1}
  = M^{-1}K\bigl(F^{-1} - M^{-1}\bigr)
  = M^{-1}KF^{-1}KM^{-1} \succeq 0 ,
\]
the last matrix being a congruence of $F^{-1} \succ 0$. Hence
$\Omega_K \succeq I_G - UF^{-1}U^{\top}$: the penalized covariance is \emph{larger} than
the maximum likelihood one, so a test referred to the classical covariance understates
the residual variance and is anti-conservative on that account alone.
Table~\ref{tab:geometry} confirms it empirically --- $\tr(\Omega_K)$ exceeds
$\tr(\Omega_{\mathrm{MLE}})$ in every row.

The covariance therefore pushes in the \emph{same} direction as the non-centrality rather
than against it. It is nevertheless the much smaller of the two effects --- the trace gap
is $0.23$ in design~A and $0.92$ in design~B, against a size distortion that reaches
$1.000$ --- and Proposition~\ref{prop:cancel} disposes of it entirely, since after
$\hat\bmu$ is subtracted the covariance is exactly the maximum likelihood one. The
non-centrality $\bmu_K$ is what has to be removed by hand.
\end{remark}

\subsection{Proof of Proposition 2}\label{sec:proof2}

\begin{proposition}\label{prop:cancel}
Let $\bbt = \bbh + F^{-1}K\bbh$ and $\hat\bmu = UM^{-1}K\bbt$. Then, to first order,
$\mathbb{E}(\br - \hat\bmu) = \boldsymbol{0}$ and
$\var(\br - \hat\bmu) = I_G - UF^{-1}U^{\top}$.
\end{proposition}

\begin{proof}
The key is that $\bbt$ is exactly unbiased to first order. Since
$I + F^{-1}K = F^{-1}(F+K) = F^{-1}M$,
\[
  \bbt = (I + F^{-1}K)\bbh = F^{-1}M\bbh
       = F^{-1}M\bigl[\bb + M^{-1}\{X^{\top}(\by-\bpi) - K\bb\}\bigr] + o_p(1),
\]
and expanding, $F^{-1}M\bb = F^{-1}(F+K)\bb = \bb + F^{-1}K\bb$, so the $F^{-1}K\bb$
terms cancel and
\begin{equation}\label{eq:btilde}
  \bbt = \bb + F^{-1}X^{\top}(\by-\bpi) + o_p(1).
\end{equation}
That is, the one-step correction returns precisely the first-order expansion of the
unpenalized maximum likelihood estimator, whether or not that estimator exists.

Hence $\hat\bmu = UM^{-1}K\bb + UM^{-1}KF^{-1}X^{\top}(\by-\bpi) + o_p(1)$, and
subtracting from \eqref{eq:rlinear} the non-random terms cancel exactly:
\begin{align*}
  \br - \hat\bmu &= \bigl(A - UM^{-1}KF^{-1}X^{\top}\bigr)(\by - \bpi) + o_p(1), \\
  A - UM^{-1}KF^{-1}X^{\top} &= V^{-1/2}C - UM^{-1}(I + KF^{-1})X^{\top}.
\end{align*}
Now the cancellation that gives the proposition its content:
\begin{equation}\label{eq:cancel}
  M^{-1}(I + KF^{-1}) = M^{-1}(F+K)F^{-1} = M^{-1}MF^{-1} = F^{-1}.
\end{equation}
So $\br - \hat\bmu = A^{*}(\by-\bpi) + o_p(1)$ with
$A^{*} = V^{-1/2}C - UF^{-1}X^{\top}$, whose mean is zero and whose variance is
\[
  A^{*}WA^{*\top} = I_G - UF^{-1}U^{\top} - UF^{-1}U^{\top} + UF^{-1}FF^{-1}U^{\top}
                  = I_G - UF^{-1}U^{\top}. \qedhere
\]
\end{proof}

\begin{remark}
The estimation noise that the plug-in $\hat\bmu$ injects is exactly the amount needed to
convert the penalized covariance back into the maximum likelihood covariance. This is a
cancellation, not an approximation: \eqref{eq:cancel} is an algebraic identity. It is the
reason the practitioner does not have to learn a new reference distribution.

It also identifies the one quantity the procedure depends on badly, namely $F^{-1}$. When
$F$ is near-singular --- light penalty with $p/n$ large --- the correction
$F^{-1}K\bbh$ can overshoot. Section~\ref{sec:overshoot} measures this.
\end{remark}

\subsection{Proof of Proposition 3}\label{sec:proof3}

\begin{proposition}\label{prop:atten}
Let $(\eta^{\ast},\hat\eta^{\ast})$ be standard bivariate normal with correlation $\rho$.
Then $\mathbb{E}\{\mathrm{He}_k(\eta^{\ast}) \mid \hat\eta^{\ast}\}
= \rho^{k}\,\mathrm{He}_k(\hat\eta^{\ast})$, and consequently a departure equal to the
degree-$k$ Hermite component of the true linear predictor contributes non-centrality
attenuated by $\rho^{2k}$ to a test that groups on the fitted predictor.
\end{proposition}

\begin{proof}
Use the generating function $\sum_{k\ge0} (t^k/k!)\,\mathrm{He}_k(x) = \exp(tx - t^2/2)$.
Given $\hat\eta^{\ast} = v$, we have $\eta^{\ast} \sim \mathcal{N}(\rho v, 1-\rho^{2})$, so
\begin{align*}
  \mathbb{E}\bigl\{\exp(t\eta^{\ast} - t^{2}/2) \mid \hat\eta^{\ast} = v\bigr\}
  &= \exp(-t^{2}/2)\exp\bigl\{t\rho v + \tfrac12 t^{2}(1-\rho^{2})\bigr\} \\
  &= \exp\bigl\{(t\rho)v - (t\rho)^{2}/2\bigr\}
   = \sum_{k\ge0}\frac{(t\rho)^{k}}{k!}\mathrm{He}_k(v).
\end{align*}
Comparing coefficients of $t^{k}$ gives the conditional expectation. For the second
claim, a grouped test on $\hat\eta^{\ast}$ is a function of the data only through group
means, so the departure it responds to is
$\mathbb{E}\{g(\eta^{\ast}) \mid \hat\eta^{\ast}\}$. Taking
$g = \mathrm{He}_k$ gives $\rho^{k}\mathrm{He}_k(\hat\eta^{\ast})$; the non-centrality of a
quadratic-form statistic is proportional to the squared norm of that mean vector, hence
to $\rho^{2k}$.
\end{proof}

\begin{remark}\label{rem:tauabsorbs}
The visible effect size $\tau$ of the main text has
$\sd[\mathbb{E}\{g(\boldsymbol{x})\mid\hat\eta\}]$ in its numerator, which already
contains the factor $\rho^{k}$. So $\tau$ absorbs the attenuation by construction, and no
residual dimensional offset should be expected on a $\tau$ axis --- as
the index-aligned panel of the power figure in the main text confirms. Proposition~\ref{prop:atten} governs how much
$\tau$ a given nominal $\gamma$ buys, not the shape of power as a function of $\tau$.
\end{remark}


% =====================================================================
\section{Alternative corrections that fail}\label{sec:failures}
% =====================================================================

% MANIFEST section 4 -- GREEN throughout except the last row.
The prepivoted procedure of the main text has four neighbouring constructions that are
simpler and that do not work. We report them because each failure is diagnostic, and
because a reader's first reaction to Proposition~\ref{prop:cancel} is reasonably ``why
not just subtract $\hat\bmu$ and use the $\chi^2$ reference?''.

\begin{table}[htbp]
\centering
\caption{Rejection rate under a correctly specified model, nominal level $0.05$, for the
constructions that fail. Design~B entries are given at the penalties stated. Every row
was run on the same designs and groupings as the main results.}
\label{tab:failures}
\begin{tabular}{p{6.0cm}p{5.0cm}p{2.7cm}}
\toprule
Construction & Rejection rate under $H_0$ & Verdict \\
\midrule
(a) Analytic corrected statistic referred to $\Omega_{\mathrm{MLE}}$
  & A: $0.036$--$0.064$; B: $0.050 / 0.026 / 0.016$ at $\lambda = 20/50/416$
  & conservative in B \\
(b) Second-order Taylor expansion of $\bmu$
  & A: $0.043$--$0.057$; B: $0.240 / 0.107 / 0.103$
  & overshoots by $\approx 2\times$ \\
(c) Debiased bootstrap with the \emph{plain} (uncorrected) statistic
  & $0.000$--$0.060$
  & conservative everywhere \\
(d) Prepivoting with the generator $\hat\bpi_K$ from the \emph{penalized} fit
  & $0.008 / 0.025$
  & conservative \\
\midrule
\textbf{(e) Prepivoting with generator $\bpi(\bbt)$ and the corrected statistic in
  \emph{both} worlds}
  & see the corrected-size table of the main text
  & \textbf{valid} \\
\midrule
(f) Wild Rademacher multiplier bootstrap of $\br - \hat\bmu$, refit-free
  & $0.058 / 0.033$
  & promising; one run only \\
\bottomrule
\end{tabular}
\end{table}

Three comments on Table~\ref{tab:failures}.

(a) The analytic correction is enough in design~A and not in design~B. This is the
boundary of the first-order theory, and it is the reason the paper does not stop at
Proposition~\ref{prop:cancel}.

(b) It is natural to respond to (a) by taking the expansion one term further. This makes
matters worse, not better, and the diagnostic shows why: at $\lambda = 416$ the
first-order $\bmu_1$ tracks the empirical mean of the grouped residuals closely --- the
empirical mean ranges over $(-3.79, 3.65)$ and $\bmu_1$ over $(-3.41, 3.36)$ --- whereas
the second-order $\bmu_2$ ranges over $(-6.32, 6.53)$, roughly twice too large. The
expansion is at its radius of convergence, so adding terms diverges rather than refines.
This is the clearest single piece of evidence that the problem at $p/n = 0.25$ is not one
that a better expansion solves.

(d) The generator matters as much as the statistic. Generating from the penalized fit
builds the shrinkage into the null world, so the bootstrap statistics carry the same
displacement as the observed one and the test loses its ability to reject. Generating
from $\bpi(\bbt)$ instead is what makes the null world the \emph{structural} null of
$H_0^{\mathrm{str}}$ of Section~2.2 of the main text, which is the hypothesis we mean to
test, rather than the calibration null $H_0^{\mathrm{cal}}$, which the penalized fit
violates by construction.

Row (f) is a genuinely attractive alternative because it avoids refitting entirely, and
therefore costs a fraction of the computation. We report it as promising rather than
established: it rests on a single run at $B = 120$, which gives a Monte Carlo standard
error of about $0.02$, and we have not examined its power. It is the natural subject of
follow-up work.


% =====================================================================
\section{Simulation designs}\label{sec:designs}
% =====================================================================

\textbf{Design~A} ($n = 500$, $p = 5$). Covariates are independent standard normal. The
linear predictor is $\eta_i = \boldsymbol{x}_i^{\top}\bb_0$ with
$\bb_0 = (0.5, -0.4, 0.3, -0.3, 0.2)^{\top}$ and no intercept, giving
$\sd(\eta) \approx 0.79$ and a balanced response. Penalties examined are
$\lambda \in \{100, 137, 200\}$, equivalently
$\lambda_{\mathrm{g}} \in \{0.20, 0.274, 0.40\}$.

\textbf{Design~B} ($n = 400$, $p = 100$, $p/n = 0.25$). Covariates are multivariate
normal with an autoregressive correlation $\Sigma_{jk} = 0.7^{|j-k|}$. The coefficient
vector cycles $(0.35, -0.30, 0.25, -0.20, 0.15)$ over the $100$ coordinates and is scaled
by $0.8887$ so that $\|\bb_0\| = 2.31$. Penalties examined are
$\lambda \in \{50, 416, 1000\}$, equivalently
$\lambda_{\mathrm{g}} \in \{0.125, 1.04, 2.5\}$.

In both designs the fitting is ridge logistic regression by iteratively reweighted least
squares on the penalized normal equations, with the intercept unpenalized, to a tolerance
of $10^{-9}$ and at most $60$ iterations. The penalty scale identity
$\lambda = n\lambda_{\mathrm{g}}$ against \texttt{glmnet} was verified numerically to
$1.6\times10^{-9}$.

Grouping is by equal-frequency deciles of the fitted probability throughout unless stated
otherwise. The EDGE basis is the orthogonal cubic polynomial basis in the group-mean
fitted probability, giving a three-dimensional quadratic form; the implementation used
here was verified against \texttt{ebrahim.gof::edge.gof} to a difference of
$0.00\times10^{0}$ on three datasets.

Under the null, data are generated from the design's own logistic model, so every
rejection reported as ``size'' is a false rejection of a model that is exactly correct.


% =====================================================================
\section{Full grids and sensitivity}\label{sec:grids}
% =====================================================================

\subsection{Size of the uncorrected tests}\label{sec:naivegrid}

% MANIFEST section 2 -- GREEN. This is the exact grid drawn as the penalty-path figure
% of the main text; it lives here so the figure does not have to carry five digits.
\begin{table}[htbp]
\centering
\caption{Rejection rate of the uncorrected tests under a \emph{correctly specified}
model, at nominal level $0.05$, over the ridge penalty. This is the grid plotted in
panels~A and~B of the penalty-path figure of the main text. Design~A has $n = 500$ and
$p = 5$; design~B has $n = 400$ and $p = 100$ with an autoregressive covariate
correlation of $0.7$. Penalties are on the \citet{friedman2010} scale
$\lambda_{\mathrm{g}} = \lambda/n$. Based on $500$ replications, Monte Carlo standard
error at most $0.022$. The design~B mean absolute error row is the quantity plotted on
the horizontal axis of panel~C.}
\label{tab:naivegrid}
\begin{tabular}{lcccccc}
\toprule
Design A: $\lambda_{\mathrm{g}}$ & 0 & 0.02 & 0.05 & 0.10 & 0.20 & 0.274 \\
\midrule
Hosmer--Lemeshow & 0.068 & 0.068 & 0.118 & 0.362 & 0.916 & 0.988 \\
EDGE basis       & 0.034 & 0.046 & 0.170 & 0.888 & 1.000 & 1.000 \\
\midrule
Design B: $\lambda_{\mathrm{g}}$ & 0 & 0.01 & 0.05 & 0.124 & 0.513 & 1.5 \\
\midrule
Hosmer--Lemeshow & 0.178 & 0.432 & 0.990 & 0.994 & 1.000 & 1.000 \\
EDGE basis       & 0.284 & 0.994 & 1.000 & 1.000 & 1.000 & 1.000 \\
$\mathrm{MAE}(\hat{\bpi},\bpi_0)$ & 0.1859 & 0.1469 & \textbf{0.1330} & 0.1336 & 0.1606 & 0.1840 \\
\bottomrule
\end{tabular}
\end{table}

Table~\ref{tab:naivegrid} gives the headline grid. Two features of it are worth recording
here. Over $500$ replications of both designs, the one-standard-error cross-validation
rule selected $\lambda_{\mathrm{g}} > 0.05$ in \emph{every} replication, so the columns at
which the uncorrected test has already failed are the columns practice actually visits.
And in design~B the unpenalized column is itself invalid, at $0.178$ and $0.284$; this is
the \citet{sur2019} and \citet{candes2020} phenomenon and it means that declining to
penalize is not an escape.

\subsection{The effect of the correction on the residual geometry}\label{sec:geometry}

% MANIFEST 8h(a) -- GREEN (deterministic given the design).
\begin{table}[htbp]
\centering
\caption{Effective degrees of freedom of the grouped residual vector, $G = 10$. ``Removed''
is $\mathrm{tr}(UF^{-1}U^{\top})$, the number of decile directions absorbed by estimation.
``EDGE-3'' is the effective dimension the three EDGE directions actually receive out of a
nominal three. Averaged over $60$ replications; these are geometric quantities and vary
little between replications.}
\label{tab:geometry}
\begin{tabular}{llccccc}
\toprule
Design & $\lambda$ & $\mathrm{df}(\Omega_K)$ & $\mathrm{df}(\Omega_{\mathrm{MLE}})$
& Removed & EDGE-3 & $\max\mathrm{eig}(\Omega_{\mathrm{MLE}})$ \\
\midrule
A & 100  & 8.20 & 7.97 & 2.03 & 2.020 & 1.000 \\
A & 137  & 8.27 & 7.97 & 2.03 & 2.021 & 1.000 \\
A & 200  & 8.37 & 7.97 & 2.03 & 2.021 & 1.000 \\
\midrule
B & 50   & 6.89 & 5.97 & 4.03 & 1.568 & 0.853 \\
B & 416  & 7.80 & 6.03 & 3.97 & 1.546 & 0.845 \\
B & 1000 & 8.19 & 6.04 & 3.96 & 1.539 & 0.844 \\
\bottomrule
\end{tabular}
\end{table}

Table~\ref{tab:geometry} quantifies what dimension costs a grouped test, independently of
any particular alternative.

(a) Estimation absorbs $2.03$ of the ten decile directions in design~A and $4.03$ in
design~B --- almost exactly twice as many. A test with ten groups is effectively a
six-dimensional test at $p/n = 0.25$.

(b) In design~A the largest eigenvalue of $\Omega_{\mathrm{MLE}}$ is exactly $1.000$:
some direction of the grouped residual vector is left entirely untouched by estimation.
In design~B the largest eigenvalue never exceeds $0.853$. No direction survives intact,
which is a compact way of saying what $p/n$ costs.

(c) The three EDGE directions receive only about $1.55$ effective dimensions in design~B
against $2.02$ in design~A.

\subsection{Behaviour of the debiasing at light penalty}\label{sec:overshoot}

% MANIFEST 8h(b) -- GREEN for the norms.
Proposition~\ref{prop:cancel} depends on $F^{-1}$, which is poorly conditioned when the
penalty is light and $p/n$ is large. Table~\ref{tab:overshoot} measures the consequence.

\begin{table}[htbp]
\centering
\caption{Norm of the debiased estimate relative to the truth, averaged over $60$
replications. The debiasing recovers the scale of $\bb_0$ almost exactly except at the
lightest penalty in design~B, where it overshoots by $26\%$.}
\label{tab:overshoot}
\begin{tabular}{llccc}
\toprule
Design & $\lambda$ & $\|\bb_0\|$ & $\|\bbh\|$ & $\|\bbt\|/\|\bb_0\|$ \\
\midrule
A & 100  & 0.79 & 0.42 & 0.99 \\
A & 137  & 0.79 & 0.36 & 0.98 \\
A & 200  & 0.79 & 0.29 & 0.96 \\
\midrule
B & 50   & 2.31 & 0.84 & \textbf{1.26} \\
B & 416  & 2.31 & 0.27 & 1.05 \\
B & 1000 & 2.31 & 0.17 & 1.02 \\
\bottomrule
\end{tabular}
\end{table}

The overshoot is confined to design~B at $\lambda = 50$ and is within $5\%$ of the truth
in all five other cells, which is what the conditioning argument predicts: heavier
penalties make $M = F + K$ better conditioned, so the correction is stable again.

We tested whether the overshoot is consequential by intervening on it directly, which is
a stronger design than correlating it with anything. Holding the statistic, the grouping
and the observed data fixed, we replaced the generator $\bpi(\bbt)$ by
$\bpi\{\bbh + s(\bbt - \bbh)\}$ with $s = 0.71$, chosen so that the generator has the
same norm as $\bb_0$; the achieved ratio moved from $1.26$ to $1.00$. Over $400$
replications at $\lambda = 50$ the rejection rate on the EDGE basis was $0.0500$ before
the intervention and $0.0550$ after it, a difference well inside the Monte Carlo standard
error of $0.011$, and the ratio of bootstrap to observed statistic variance moved from
$0.72$ to $0.74$. Removing the overshoot does not remove the conservatism.

Two cautions about how far this intervention settles the question, because it does not
settle it as cleanly as the arithmetic above suggests. First, the pre-intervention arm
returned $0.0500$, whereas the pooled size of that same cell over $4000$ replications is
$0.0370$; the $400$-replication block used here did not itself display the conservatism
the intervention was meant to remove. Second, the quoted $0.011$ is an unpaired Monte
Carlo standard error, whereas the design is paired --- the statistic, the grouping and the
data are held fixed --- so the experiment's true resolution is finer than $0.011$ and we
have not measured it. What can be said without either caveat is that the bootstrap
reference is \emph{less} dispersed than the observed statistic, $\var(S^{*})/\var(S) = 0.72$,
which pushes towards over-rejection and so cannot account for under-rejection either.

We therefore record the overshoot as an unexplained correlate rather than a refuted
hypothesis. It is large, it is confined to exactly the anomalous cell, and it has a clean
mechanism in the conditioning of $F$; what we can show is that removing it does not remove
the conservatism. The overshoot is a real caution about the
behaviour of $\bbt$ at light penalty when $p/n$ is large, and it is worth knowing for
anyone who wants to use $\bbt$ for its own sake, but it does not affect the size of the
test and we do not offer it as an account of Section~\ref{sec:contested}'s residual
conservatism, which remains unexplained.

\subsection{Size-adjusted power}\label{sec:sizeadj}

% MANIFEST 6.3 -- GREEN.
Because the EDGE basis under-rejects in design~B at light penalty
(Section~\ref{sec:contested}), a comparison at the nominal level is not quite a
comparison of like with like. We therefore repeated every power comparison at a
size-adjusted level: the $\gamma = 0$ cell of each panel supplies the null distribution
of the $p$-value, and the alternative cells are thresholded at its fifth percentile.
The thresholds were $0.0533$ in design~A on both bases, and $0.0400$ (decile) and
$0.0533$ (EDGE) in design~B.

\begin{table}[htbp]
\centering
\caption{Nominal and size-adjusted power, $2000$ replications per cell. The adjustment
moves no entry by more than $0.02$ and reverses no comparison, so the conclusions of
the power table of the main text do not depend on it.}
\label{tab:sizeadj}
\begin{tabular}{lllcccc}
\toprule
& & & \multicolumn{2}{c}{Nominal} & \multicolumn{2}{c}{Size-adjusted} \\
\cmidrule(lr){4-5}\cmidrule(lr){6-7}
Design & Departure & $\gamma$ & decile & EDGE & decile & EDGE \\
\midrule
A & ---         & 0.00 & 0.0445 & 0.0460 & 0.0520 & 0.0530 \\
A & quadratic   & 0.25 & 0.0760 & 0.1150 & 0.0820 & 0.1305 \\
A & quadratic   & 0.50 & 0.1475 & 0.2615 & 0.1660 & 0.2820 \\
A & quadratic   & 0.75 & 0.2015 & 0.3755 & 0.2175 & 0.3940 \\
A & quadratic   & 1.00 & 0.2700 & 0.4150 & 0.2885 & 0.4300 \\
A & interaction & 0.25 & 0.0595 & 0.0750 & 0.0680 & 0.0830 \\
A & interaction & 0.50 & 0.1055 & 0.1885 & 0.1185 & 0.2045 \\
A & interaction & 0.75 & 0.1940 & 0.3555 & 0.2140 & 0.3745 \\
A & interaction & 1.00 & 0.2835 & 0.4615 & 0.3080 & 0.4855 \\
\midrule
B & ---         & 0.00 & 0.0640 & 0.0495 & 0.0560 & 0.0580 \\
B & quadratic   & 0.50 & 0.0595 & 0.0485 & 0.0535 & 0.0570 \\
B & quadratic   & 1.00 & 0.0635 & 0.0600 & 0.0525 & 0.0670 \\
B & quadratic   & 1.50 & 0.0740 & 0.0630 & 0.0655 & 0.0675 \\
\bottomrule
\end{tabular}
\end{table}

\subsection{Computational cost of the procedure and of prepivoting}\label{sec:cost}

% MANIFEST sections 8e and 7.6 -- GREEN.
The obvious objection to a procedure that refits inside a bootstrap loop is its price, so
we state it. On the machine and worker count given in Table~\ref{tab:repro}, one
prepivoted test at $N_B = 149$ takes about $0.26$ seconds in design~A and about $3.3$
seconds in design~B, the difference being the cost of $149$ ridge refits at $p = 100$.
Two independent timings agree to within $2\%$. For a single analysis --- which is what a
practitioner runs --- the cost is seconds, and the simulation studies in this paper are
expensive only because they repeat that analysis thousands of times.

Prepivoting itself is not the expensive part, and it is not a statistical tax either. In
design~A at $\lambda = 100$, referring the corrected statistic to the analytic
$\Omega_{\mathrm{MLE}}$ reference gives rejection rates $0.023$, $0.147$ and $0.227$ at
$\gamma = 0$, $0.5$ and $1.0$, while prepivoting the same statistic gives $0.040$,
$0.163$ and $0.250$. Prepivoting is closer to nominal under the null \emph{and}
higher under both alternatives. This is a smaller run than the size study --- $300$
replications, so a Monte Carlo standard error of about $0.029$ --- and we would not read
the individual gaps as established; the point is only that resampling does not buy
validity at the price of power here.

\subsection{Re-estimation of the contested cell}\label{sec:contested}

% MANIFEST 8i -- GREEN.
The EDGE basis in design~B at $\lambda = 50$ is the one cell of the corrected-size table of the main text
that we estimated more than once, and we set out the sequence because it bears on how
such tables should be read.

Our first estimate, from a single block of $1000$ replications, was $0.0270$ with a
standard error of $0.0051$: four and a half standard errors below nominal, and outside
the range $[0.03, 0.08]$ that we had fixed in advance as acceptable. Before writing about
it we re-estimated the cell on two further blocks of $1500$, disjoint from the first and
from each other. Table~\ref{tab:contested} gives all three.

\begin{table}[htbp]
\centering
\caption{The contested cell, design~B at $\lambda = 50$, estimated on three independent
replication blocks. The decile basis was never in doubt. On the EDGE basis the first
block is the lowest of the three; the three are homogeneous and the pooled estimate lies
inside the acceptable range.}
\label{tab:contested}
\begin{tabular}{lccccc}
\toprule
Basis & Block 1 ($1000$) & Block 2 ($1500$) & Block 3 ($1500$)
& Homogeneity & Pooled ($4000$) \\
\midrule
Decile & 0.0630 & 0.0587 & 0.0580 & $p = 0.86$ & 0.0595 (0.0037) \\
EDGE   & 0.0270 & 0.0440 & 0.0367 & $p = 0.09$ & 0.0370 (0.0030) \\
\bottomrule
\end{tabular}
\end{table}

Three conclusions. The conservatism on the EDGE basis is \emph{real}: pooled over $4000$
replications the size is $0.0370$, with an exact $95\%$ interval of $[0.0314, 0.0433]$
and $z = -4.36$ against nominal, so it is not the kind of blemish that disappears with
more computation. Its \emph{magnitude} was nonetheless overstated by the first block by
about a third. And the pooled value lies inside $[0.03, 0.08]$, so the qualitative
verdict on the cell changes even though the direction does not.

We draw the obvious methodological moral. A rejection rate near $0.03$ estimated on
$1000$ replications has a standard error of $0.005$, so blocks will disagree by more than
one might expect, and the residual heterogeneity here ($p = 0.09$) suggests $1000$ is
thin for a tail probability this small. Any cell that falls outside a pre-specified range
should be re-estimated on an independent block before conclusions are drawn from it.

\subsection{Sensitivity to the number of groups and to the number of bootstrap
replicates}\label{sec:sensitivity}

On the glaucoma data the corrected $p$-value varies from $0.026$ at $G = 10$ to $0.172$
at $G = 20$ on the decile basis. This sensitivity is inherited from the
Hosmer--Lemeshow family and the correction does not remove it; the full $\lambda\times G$
grid is the full $\lambda\times G$ glaucoma table of the main text, and we report all twelve cells rather than a selected
one. On the same data, repeating the cross-validated cell at $G=10$ with $N_B = 999$
under different seeds gave corrected $p$-values in the range $0.015$--$0.026$ on the
decile basis and $0.034$--$0.049$ on the EDGE basis, so the reported figures are stable
to the bootstrap size.


% =====================================================================
\section{The glaucoma analysis in detail}\label{sec:glaucoma}
% =====================================================================

The \texttt{GlaucomaM} data are distributed in the R package \texttt{TH.data}. There are
$n = 196$ eyes, $98$ glaucomatous and $98$ normal, matched by age and sex, with $p = 62$
morphometric variables derived from confocal laser scanning tomography of the optic nerve
head. All variables were centred and scaled before fitting.

Design facts: $\mathrm{cond}(X^{\top}X) = 4.97\times10^{7}$, largest absolute pairwise
correlation $0.9961$, and $p/n = 0.316$. The unpenalized logistic fit does not converge;
the maximum likelihood estimate does not exist. Ten-fold cross-validation gives
$\hat\lambda_{\min} = 19.9$ and $\hat\lambda_{1\mathrm{se}} = 295.3$ on the theory scale,
the latter corresponding to $\lambda_{\mathrm{g}} = 1.51$. At
$\hat\lambda_{1\mathrm{se}}$ the fit has area under the curve $0.905$,
$\sd(\hat\pi) = 0.224$ and fitted probabilities spanning $[0.016, 0.911]$; at
$\hat\lambda_{\min}$ the area under the curve is $0.942$. The apparent calibration slope
at $\hat\lambda_{1\mathrm{se}}$ is $2.434$.

\subsection{A second dataset}\label{sec:mvf}

% MANIFEST section 8b-ter -- GREEN. Full 12-cell grid, NB = 499.
The phenomenon is not confined to one dataset. \texttt{GlaucomaMVF} (package
\texttt{ipred}) is a second confocal-tomography series, $n = 170$ eyes and $p = 63$
morphometric variables, so $p/n = 0.371$; here too the maximum likelihood estimate does
not exist. Cross-validation gives $\hat\lambda_{\min} = 10.2$ and
$\hat\lambda_{1\mathrm{se}} = 54.6$ on the theory scale.
Table~\ref{tab:mvf} reports the same grid we ran for \texttt{GlaucomaM}.

\begin{table}[htbp]
\centering
\caption{The uncorrected and corrected tests on \texttt{GlaucomaMVF}, four penalties by
three group counts, $499$ bootstrap replicates. The reversal is complete: the uncorrected
EDGE test rejects in all twelve configurations, and the corrected test rejects in none of
them on either basis.}
\label{tab:mvf}
\small
\setlength{\tabcolsep}{4pt}
\begin{tabular}{lccccccc}
\toprule
& \multicolumn{3}{c}{Uncorrected (dec / EDGE)} & & \multicolumn{3}{c}{Corrected (dec / EDGE)} \\
\cmidrule(lr){2-4}\cmidrule(lr){6-8}
$\lambda$ & $G=5$ & $G=10$ & $G=20$ & & $G=5$ & $G=10$ & $G=20$ \\
\midrule
$\hat\lambda_{\min}$ & 0.077 / 0.041 & 0.023 / 0.003 & \textbf{0.218} / 0.005 & & 0.458 / 0.318 & 0.556 / 0.350 & 0.476 / 0.298 \\
$\tfrac12\hat\lambda_{1\mathrm{se}}$ & 0.007 / 0.013 & 0.029 / 0.003 & \textbf{0.335} / 0.003 & & 0.224 / 0.826 & 0.706 / 0.552 & 0.808 / 0.636 \\
$\hat\lambda_{1\mathrm{se}}$ & 0.000 / 0.000 & 0.016 / 0.001 & \textbf{0.327} / 0.001 & & 0.688 / 0.874 & 0.876 / 0.470 & 0.998 / 0.428 \\
$2\hat\lambda_{1\mathrm{se}}$ & 0.000 / 0.000 & 0.001 / 0.000 & 0.027 / 0.000 & & 0.430 / 0.852 & 0.618 / 0.564 & 0.840 / 0.426 \\
\bottomrule
\end{tabular}
\end{table}

Three things should be said about Table~\ref{tab:mvf}.

(a) The reversal is cleaner here than on \texttt{GlaucomaM}. The corrected test returns
$p$ between $0.224$ and $0.998$ on the decile basis and between $0.298$ and $0.874$ on
the EDGE basis, rejecting in none of the twelve cells, while the uncorrected EDGE test
rejects in all twelve. Repeating the cross-validated cell at $G = 10$ with $N_B = 999$
under five seeds gives corrected $p$-values in $[0.883, 0.916]$ on the decile basis and
$[0.419, 0.465]$ on the EDGE basis, so the conclusion is not an artefact of the bootstrap
size.

(b) The uncorrected \emph{decile} test does not reject in four of the twelve cells,
shown in bold, three of them at $G = 20$. We therefore do not claim that the uncorrected
test rejects everywhere on this dataset; the claim holds on the EDGE basis, and on the
decile basis at the finer groupings.

(c) We report this dataset in the supplement rather than the main text because it adds
corroboration rather than a new argument. A reader who doubts that \texttt{GlaucomaM} is
representative should read it as the reply.

We also examined \texttt{Sonar} (package \texttt{mlbench}; $n = 208$, $p = 60$,
$p/n = 0.288$, maximum likelihood separates) and observed the same qualitative reversal,
but only at a single penalty and group count. We do not report its $p$-values, because a
single configuration cannot be distinguished from a favourable choice of one.

We also examined \texttt{wpbc} and rejected it as an example. At the cross-validated
penalty the largest fitted coefficient is $0.00099$ and the fitted probabilities span
$[0.234, 0.242]$: cross-validation had selected an essentially intercept-only model, so
its apparent clean reversal from $p = 0.0008$ to $p = 0.80$ is an artefact of having no
model to test rather than a demonstration of the method. We report this because it is the
kind of example that would flatter the method if reported uncritically.


% =====================================================================
\section{Reproducibility}\label{sec:repro}
% =====================================================================

All computation is in R. Penalized fits use the iteratively reweighted least squares
routine described in Section~\ref{sec:designs}; \texttt{glmnet} \citep{friedman2010} is
used for cross-validation and for the penalty-scale check. Parallel runs use PSOCK
clusters with \texttt{parallel::clusterSetRNGStream} for reproducible streams; the stream
seed is given with each result below.

\begin{table}[htbp]
\centering
\caption{Scripts, seeds and cost for every reported experiment.}
\label{tab:repro}
\small
\setlength{\tabcolsep}{4pt}
\begin{tabular}{p{4.0cm}p{4.3cm}p{4.0cm}p{2.2cm}}
\toprule
Result & Script & Seed / replication & Cost \\
\midrule
Proposition 1 numerical check & \texttt{corrected\_test.R} & fixed design & seconds \\
Uncorrected size grid (Table 1) & \texttt{lambda\_sweep.R}, \texttt{hd\_sweep.R}
  & $B = 500$ & minutes \\
Corrected size (Table 2) & \texttt{T23\_rerun.R}
  & stream $20260807$, seeds $1001$--$2000$ & $11.2$ min, 22 workers \\
Contested cell (Table~\ref{tab:contested}) & \texttt{verify\_E13.R}
  & seeds $1$--$1500$ and $100001$--$101500$ & $15$ min, 22 workers \\
Power (Table 3, Table~\ref{tab:sizeadj}) & \texttt{T25\_power\_highB.R}
  & stream $20260810$, $B = 2000$ & $37$ min, 22 workers \\
Visible effect size (Figure 1) & \texttt{T22\_tau.R}, \texttt{plot\_tau.R},
  \texttt{R/fig1\_power\_ek.R}
  & known-null MC, $B = 2000$ null, $B_1 = 1000$ & minutes \\
Oracle grouping (Table 5) & \texttt{oracle\_knownnull.R}
  & $B_0 = 2000$, $B_1 = 1000$, both arms exact & minutes \\
Residual geometry (Table~\ref{tab:geometry}) & \texttt{omega\_spectrum.R}
  & stream $20260808$, $60$ reps & $26$ s, 22 workers \\
Overshoot intervention (\S\ref{sec:overshoot}) & \texttt{overshoot\_test.R}
  & stream $20260809$, $400$ reps & $244$ s, 22 workers \\
Failed repairs (Table~\ref{tab:failures}) & \texttt{round2.R}, \texttt{cand\_C.R}
  & $B = 120$--$300$ & minutes \\
Glaucoma analysis & \texttt{glaucoma\_deep.R}, \texttt{R/fig3\_calib\_ek.R}
  & \texttt{set.seed(1)}, $N_B = 499$ & minutes \\
Second dataset (\S\ref{sec:mvf}) & \texttt{T26\_app.R}, \texttt{extract\_T26b.R}
  & \texttt{set.seed(1)}, $N_B = 499$ & minutes \\
Cost, and the price of prepivoting & \texttt{prepivot\_cost.R}
  & $B = 300$, design~A $\lambda = 100$ & minutes \\
Conditioning on the application & \texttt{glaucoma\_conditioning.R}
  & \texttt{set.seed(1)} & seconds \\
Generator diagnostic & \texttt{check\_generator.R}
  & \texttt{set.seed(1)} & seconds \\
Covariance ordering & \texttt{check\_omega\_sign.R}
  & $500$ random $(F,K)$ draws & seconds \\
\bottomrule
\end{tabular}
\end{table}

Replication-level $p$-values are archived for the power study
(\texttt{T25\_power\_pvalues.csv}) and for the head-to-head comparison with the
high-dimensional alternatives (\texttt{T24\_rivals\_pvalues.csv}), so that every cell of
the power tables and every size-adjusted threshold can be re-derived and its Monte Carlo
standard error recomputed independently. The remaining experiments archive
per-replication rejection indicators and the saved fitted objects
(\texttt{*\_results.rds}), which is enough to reproduce each reported rate but not to
re-threshold it at a different level. The complete archive, including these files, the
figure scripts and \texttt{\_ek\_theme.R}, is available at
\url{https://github.com/ebrahimkhaled/gof-penalized-paper} and deposited at
\url{https://doi.org/10.5281/zenodo.21900115}.

\bibliographystyle{plainnat}
\begin{thebibliography}{9}
\bibitem[Cand\`es and Sur(2020)]{candes2020}
Cand\`es, E.J. and Sur, P. (2020). The phase transition for the existence of the maximum
likelihood estimate in high-dimensional logistic regression. \emph{The Annals of
Statistics}, \textbf{48}(1), 27--42.
\bibitem[Chernoff and Lehmann(1954)]{chernoff1954}
Chernoff, H. and Lehmann, E.L. (1954). The use of maximum likelihood estimates in
$\chi^2$ tests for goodness of fit. \emph{The Annals of Mathematical Statistics},
\textbf{25}(3), 579--586.
\bibitem[Friedman et~al.(2010)]{friedman2010}
Friedman, J., Hastie, T. and Tibshirani, R. (2010). Regularization paths for generalized
linear models via coordinate descent. \emph{Journal of Statistical Software},
\textbf{33}(1), 1--22.
\bibitem[Moore and Spruill(1975)]{moore1975}
Moore, D.S. and Spruill, M.C. (1975). Unified large-sample theory of general chi-squared
statistics for tests of fit. \emph{The Annals of Statistics}, \textbf{3}(3), 599--616.
\bibitem[Sur and Cand\`es(2019)]{sur2019}
Sur, P. and Cand\`es, E.J. (2019). A modern maximum-likelihood theory for
high-dimensional logistic regression. \emph{Proceedings of the National Academy of
Sciences}, \textbf{116}(29), 14516--14525.
\end{thebibliography}

\end{document}
